\documentclass[11pt,a4paper]{article}

\usepackage[margin=0.5in, top=3cm, bottom=2cm]{geometry}
\usepackage[spanish, activeacute]{babel}
\usepackage[utf8]{inputenc}
\usepackage{amsthm}
\usepackage{amsmath}
\usepackage{amsfonts}
\usepackage{amssymb}
\usepackage{graphicx} %Para incluir el logo de la UBA
\usepackage{caratula} %Para armar el cuadro de integrantes
\usepackage{hyperref} %Para escribir urls
\usepackage{pdfpages} %Para incluir pdf's
\usepackage{todonotes}
\usepackage{float}
\usepackage[linesnumbered]{algorithm2e}


\graphicspath{{imagenes/}}

%Cosas para escribir codigo fuente
%Fuente: http://en.wikibooks.org/wiki/LaTeX/Source_Code_Listings
\usepackage{listings}
\usepackage{color}

\setcounter{secnumdepth}{5}

\definecolor{mygreen}{rgb}{0,0.65,0.25}
\definecolor{mygray}{rgb}{0.5,0.5,0.5}
\definecolor{myorange}{rgb}{1,0.4,0.2}
\definecolor{myblue}{rgb}{0,0,0.85}
\definecolor{myred}{rgb}{0.85,0,0}

%Configuracion para los listings
\lstset{ %
  backgroundcolor=\color{white},   % choose the background color; you must add \usepackage{color} or \usepackage{xcolor}
  basicstyle=\small,        % the size of the fonts that are used for the code
  breakatwhitespace=false,         % sets if automatic breaks should only happen at whitespace
  breaklines=true,                 % sets automatic line breaking
  captionpos=b,                    % sets the caption-position to bottom
  commentstyle=\color{myred},    % comment style
  deletekeywords={...},            % if you want to delete keywords from the given language
  escapeinside={\%*}{*)},          % if you want to add LaTeX within your code
  extendedchars=true,              % lets you use non-ASCII characters; for 8-bits encodings only, does not work with UTF-8
  frame=single,                    % adds a frame around the code
  keywordstyle=\color{myblue},       % keyword style
  language=Java,                 % the language of the code
  morekeywords={*,...},            % if you want to add more keywords to the set
  numbers=left,                    % where to put the line-numbers; possible values are (none, left, right)
  numbersep=5pt,                   % how far the line-numbers are from the code
  numberstyle=\tiny\color{mygray}, % the style that is used for the line-numbers
  rulecolor=\color{black},         % if not set, the frame-color may be changed on line-breaks within not-black text (e.g. comments (green here))
  showspaces=false,                % show spaces everywhere adding particular underscores; it overrides 'showstringspaces'
  showstringspaces=false,          % underline spaces within strings only
  showtabs=false,                  % show tabs within strings adding particular underscores
  stepnumber=1,                    % the step between two line-numbers. If it's 1, each line will be numbered
  stringstyle=\color{myorange},     % string literal style
  tabsize=2,                       % sets default tabsize to 2 spaces
  title=\lstname                   % show the filename of files included with \lstinputlisting; also try caption instead of title
}

\renewcommand{\lstlistingname}{C\'{o}digo}

\lstset{language=Java,caption={},label=DescriptiveLabel}

\begin{document}

\integrante{Altman, Martín}{COMPLETAR}{COMPLETAR}
\integrante{Chapresto, Matías}{COMPLETAR}{COMPLETAR}
\integrante{Zylber, Ariel}{530/12}{arielzylber@gmail.com}

\def\Materia{Teoría de las Comunicaciones}
\def\Titulo{Trabajo Pr\'{a}ctico 1}
\def\Fecha{20 de Abril de 2016}

%----- CARATULA -----%

\thispagestyle{empty}

\begin{center}
	\includegraphics[scale = 0.25]{imagenes/logo_uba.jpg}
\end{center}

\begin{center}
	{\textbf{\large UNIVERSIDAD DE BUENOS AIRES}}\\[1.5em]
	{\textbf{\large Departamento de Computaci\'{o}n}}\\[1.5em]
    {\textbf{\large Facultad de Ciencias Exactas y Naturales}}\\
    \vspace{35mm}
    {\LARGE\textbf{\Materia}}\\[1em]    
    \vspace{15mm}
    {\Large \textbf{\Titulo}}\\[1em]
    \vspace{15mm}
    {\textbf{\Large \Fecha}}\\
    \vspace{15mm}
    \textbf{\tablaints}
\end{center}

\newpage
\thispagestyle{empty}
\tableofcontents

\parskip=5pt
\setlength{\parindent}{0pt}

\newpage
\setcounter{page}{1}
\pagenumbering{arabic}
\pagestyle{plain}

\section{Introducción}
En este trabajo práctico se propone escuchar paquetes Ethernet pasivamente de tres diferentes redes para luego realizar estadísticas sobre cierta información provista por los paquetes con el objetivo de aplicar los conceptos de teoría de la infromación visto en casos reales para estudiar el funcionamiento de ciertas redes en concreto.\\
Primero, con ayuda del script \textbf{PEDROPCAPIEDRAS} utilizamos el paquete de python conocido como Scapy para capturar los paquetes de una red y guardarlos en un archivo para poder procesarlos más tarde.\\
Luego, el script \textbf{SCAPPY DOO} se encarga de levantar los paquetes leídos en una medición para realizar los cálculos estadísticos sobre ellos. Vamos a dividir la documentación de los objetos calculados entre los que se usan para distinguir protocolos y los usados para distinguir nodos o hosts de la red.\bigskip

Con respecto a los protocolos, calculamos:
\begin{itemize}
\item La cantidad de paquetes enviados con cada protocolo. Tomamos en cuenta sólo los protocolos para los que existe algún paquete que lo utiliza. Notar que los protocolos descartados no aportan mucha información a los datos calculados posteriormente. Es decir, consideramos la fuente $S = \{ s_1, \ldots, s_n \}$ donde cada $s_i$ representa el protocolo de algún paquete de la medición.
\item A partir de esto, para cada protocolo representado por algún $s_i$, calculamos la probabilidad de que un paquete de la muestra al azar utilice ese protocolo. Si $P$ es la cantidad total de paquetes escuchados y $c_i$ la cantidad de paquetes de la muestra que usan el protocolo simbolizado por $s_i$, $p(s_i) = \frac{c_i}{P}$ es la probabilidad de que un paquete use ese protocolo.
\item Por último, calculadas todas las $p(s_i)$, calculamos por definición la entropía de la fuente dada por $$H(S) = - \sum_{i=1}^{n} p(s_i) \cdot log_2(p(s_i))$$.
\end{itemize}
\bigskip
Para distinguir los nodos de la red, proponemos como fuente $S = {s_1, \ldots, s_n}$, donde cada $s_i$ representa la dirección IP de origen de un paquete ARP. Nuevamente, como el caso anterior consideramos sólo los símbolos que corresponden a la dirección de origen de algún paquete de la muestra.\\
Además, decidimos separar el muestreo entre los paquetes ARP con operación who-is de los de operación is-at, generando dos conjuntos de muestras disjuntos. Nos pareció que de esta manera podríamos tener un analisis mas detallado e interesante de la red, distinguiendo entre nodos que preguntan y nodos que responden.\\
Decidimos tomar la dirección de origon del paquete y no la de destino ya que para los paquetes con operación who-is, la dirección de destino siempre es un broadcast a la red, por lo que no nos brindará información sobre los hosts de la red. Y para los paquetes con operación is-at, la dirección de destino será el emisor de un who-is, por lo que estaríamos contaminando la muestra de los is-at con información de los who-is que ya habremos analizado en la otra división de la muestra. (POR QUË NO ELEGIMOS DIRECCIONES MAC?)\\
Para cada uno de los dos muestreos claculamos los mismos valores calculados para el caso de los protocolos (cantidad de paquetes por símbolo, probabilidad de cada símbolo y entropía de la fuente) por lo que no volveremos a detallar el procedimiento.\bigskip

Habiendo definido todo esto procedemos a explicar cada experimento con su correspondiente resultado y conclusiones. 

\section{Experimentación}
\subsection{Starbucks}
Como primer experimento realizamos una captura en la red de fibertel zone disponible desde el Starbucks de Corrientes y Malabia aproximadamente a las 5 de la tarde un día de semana. En principio no estaba clara la cantidad de dispositivos que utilizaban la red, ya que si bien en el Starbucks se encontraban aproximadamente 20 personas,  la mayoría de ellos utilizando un dispositivo electrónico, no sabemos cuantas de ellas estaban conectadas a esta red ni si había más dispositivos fuera del local conectados a la red.\bigskip

Luego de escuchar la red durante 15 minutos (el script para escuchar paquetes permite configurar cuanto tiempo se desea medir), logramos capturar un total de COMPLETAR paquetes. En total, registramos paquetes de ¿4? protocolos distintos. Presentamos un gráfico mostrando para cada protocolo la probabilidad de que un paquete de la muestra sea de ese protocolo. Como hay grandes diferencias de magnitud entre la cantidad de paquetes del protocolo más frecuente con el resto, decidimos que era visualmente mas informativo mostrar los valores en escala logarítmica, es decir, mostrar $log(p(s_i))$ para cada protocolo $s_i$. Como las $p(s_i)$ están entre $0$ y $1$ este logarítmico siempre será negativo, por lo que nos parece más natural mostrar en su lugar $-log(p(s_i))$. Por último, como las entropía de la fuente ($H(S)$) es la esperanza de estos valores, decidimos agregarla como una columna al gráfico para poder comparar como se comporta la información que brinda cada protocolo con respecto de la media.

\end{document}
